%% ============================================================
%% 分片 06：原文第 3 节「Topological Field Theory From ${\EUN}=4$ Super Yang-Mills Theory」（原文 glshort.tex 起于第 1216 行）
%% 由 tools/harvmac2tex.py 从 arXiv 源码（hep-th/0604151v3）自动生成。
%% 只译英文散文；公式、编号、\label、\cite 一律不改。
%% ============================================================

\subsection{一族 $A$-模型}

\label{sec:familya}

我们处于如下一般情形。我们有一个量子理论，它带有费米对称性 $Q=uQ_\ell+vQ_\rr$，且满足 $Q^2=0$（模一个规范变换）。在这种情形下，路径积分一般会局域化到 $Q$-不变的场上。当 $Q$-不变场的空间 ${\eurm Y}$ 光滑、且规范群在它上面自由作用时，路径积分可以通过在 ${\eurm Y}$ 附近展开的高斯近似或单圈近似来计算。更一般地，也许需要走到更高阶，但路径积分总由 ${\eurm Y}$ 的有限阶无穷小邻域的结构决定。

只需找出 $Q$-不变的玻色场就够了，因为费米场无论如何都是无穷小的，会自动被纳入微扰论。一个玻色场在 $Q$ 下不变的条件是：对每个费米场 $\Psi$，$\delta_T\Psi$ 都为零。换句话说，${\eurm Y}$ 由把式~\eqref{eq:helpme} 与式~\eqref{eq:mimo} 的右边取零来定义。

这些表达式为零这一条件在 $u$ 与 $v$ 的标度变换下不变，因此只依赖于 $t=v/u$。式~\eqref{eq:helpme} 的右边为零给出\footnote{要得到第三个方程，我们在 $t\not=\pm i$ 时取条件 $\delta_T\eta=0$ 与 $\delta_T\tilde\eta=0$ 的一个依赖于 $t$ 的线性组合。$t=\pm i$ 时会怎样，下面再讨论。}

\begin{yizhu}[易错点]
按：脚注里"取 $\delta_T\eta=0$ 与 $\delta_T\tilde\eta=0$ 的 $t$-依赖线性组合"这句话，是理解式~\eqref{eq:loopme} 中第三条方程 $D^*\phi=0$ 来历的关键。直接由式~\eqref{eq:helpme} 令 $\delta_T\eta=0$ 与 $\delta_T\tilde\eta=0$，得到的是 $vD^*\phi+u[\bar\sigma,\sigma]=0$ 与 $-uD^*\phi+v[\bar\sigma,\sigma]=0$；把两式按适当系数组合消去 $[\bar\sigma,\sigma]$，就得到 $D^*\phi=0$，但组合系数的行列式正比于 $u^2+v^2$，故只在 $t^2\not=-1$（即 $t\not=\pm i$）时可以这样做。$t=\pm i$ 时这两式线性相关，只能得到 $D^*\phi\pm i[\bar\sigma,\sigma]=0$，这正是后文在式~\eqref{eq:toopme} 之后特别说明的情形。
\end{yizhu}
% 式 (3.29)
\setcounter{equation}{28}
\begin{equation}\label{eq:loopme}
\begin{aligned}(F-\phi\wedge\phi+tD\phi)^+ & = 0
\\
                      (F-\phi\wedge\phi-t^{-1}D\phi)^-& = 0 \\
                      D^* \phi & = 0.\\\end{aligned}
\end{equation}

对这些方程最初步的解释是取 $-t^{-1}=t$，或等价地 $t=\pm i$。对几何 Langlands 纲领很重要的正是这些 $t$ 值；此时令 ${\cal A}$ 为 ${G}_{\Bbb{C}}$ 取值的联络 ${\cal A}=A+i\phi$，其曲率为 ${\cal F}=d{\cal A}+{\cal A}\wedge{\cal A}$。于是式~\eqref{eq:loopme} 的前两个方程等价于 ${\cal F}=0$，从而 $t=\pm i$ 时式~\eqref{eq:loopme} 的解给出一个复值平坦联络，进而给出同态 $\pi_1(M)\to G_{{\Bbb{C}}}$。式~\eqref{eq:loopme} 还要求我们再添上条件 $D^*\phi=0$，并模去 $G$ 取值的规范变换。把这两步合起来，就构造出同态 $\vartheta:\pi_1(M)\to {{{G}}}_{\Bbb{C}}$ 的模空间 ${\cal Y}$（需要更明确时，我们把这个空间记作 ${\cal Y}_M$ 或 ${\cal Y}_M(G)$）。

\begin{yizhu}[含糊点拨]
按：为什么 $t=\pm i$ 时前两个方程合起来就是 ${\cal F}=0$？要点是：当 $-t^{-1}=t$ 时，$(F-\phi\wedge\phi+tD\phi)^+=0$ 与 $(F-\phi\wedge\phi+tD\phi)^-=0$ 恰好是同一个复方程 $F-\phi\wedge\phi+iD\phi=0$ 的自对偶投影与反自对偶投影；把 ${\cal A}=A+i\phi$ 的曲率展开，${\cal F}=F+iD\phi-\phi\wedge\phi$，正是这个组合。这一步是整套构造的关键：它把"扭化后的 BPS 方程"变成"平坦联络方程"，从而把四维规范场论与二维的 Higgs 丛/平坦丛语言连接起来。注意 $t=i$ 与 $t=-i$ 分别对应 ${\cal A}$ 与 $\bar{\cal A}$，具体哪个对应哪个取决于 $\phi$ 的符号约定；后文在讨论 $t$ 与 $-t$ 等价处还会回到类似的符号问题。
\end{yizhu}
这一陈述依赖于 Corlette 的一个定理 \cite{ref58}（二维情形也由 Donaldson 在 \cite{ref20} 的附录中证明）：对 $(A,\phi)$ 模去 $G$ 取值的规范变换、要求 ${\cal F}=0$ 并加上 $D^*\phi=0$，等价于直接去掉最后一个方程、只要求 ${\cal F}=0$ 以及一个稳定性条件\footnote{正如涉及 $G_{\Bbb{C}}$ 这类非紧对称群的模问题通常那样，模空间 ${\cal Y}_M(G)$ 并不字面地参数化 $\pi_1(M)\to G_{\Bbb{C}}$ 的所有共轭类。要得到好的模空间，必须去掉某些「不稳定轨道」，在这类问题里它们对应于可约化为三角形式但不能约化为直和的和同态。取 $D^*\phi=0$ 并模去 $G$ 取值的规范变换恰好去掉这些不稳定轨道，这一点在 \cite{ref58} 中得到了证明。关于稳定性所起的作用，我们在第~\ref{sec:complexstr}~节有更充分的说明。}，并模去 $G_{\Bbb{C}}$ 取值的规范变换群。

更一般地，若 $t$ 不是实数，则对取值于 $G$ 的实李代数的场 $A,\phi$，方程~\eqref{eq:loopme} 是超定的而非椭圆型的。把这些方程的厄米部分与反厄米部分分别取零，即可知对任何非实 $t$，它们都推出 ${\cal F}=0$。

然而我们想弄清楚当 $t$ 为实数时会发生什么。（对几何 Langlands 纲领最重要的 $t$ 值是 $t=\pm i$ 与 $t=\pm 1$。）此时方程是实的，我们得到一族以 $t$ 为参数的实椭圆方程。按字面写法，这些方程在 $t\not=0,\infty$ 时是正则的。不过，由式~\eqref{eq:helpme} 右边的齐次表达式可以看出，通过把第二个方程乘以 $t$、或把第一个方程乘以 $t^{-1}$，方程可以延拓到 $t=0$ 或 $t=\infty$。于是在 $t=0$ 处，我们应把第二个方程换成 $(D\phi)^-=0$；在 $t=\infty$ 处，把第一个方程换成 $(D\phi)^+=0$。这样便得到一族以 $\Bbb{RP}^1$ 为参数的实椭圆方程。

我们可以类似地理解来自式~\eqref{eq:mimo} 的未破缺超对称条件：
 % 式 (3.30)
\setcounter{equation}{29}
\begin{equation}\label{eq:toopme}
\begin{aligned}D\sigma + t[\phi,\sigma]& = 0\\
                     D\sigma - t^{-1}[\phi,\sigma]& = 0.\\
                     \end{aligned}
\end{equation}

对任何不等于 $\pm i$ 的 $t$，这些方程都推出 $D\sigma=[\phi,\sigma]=0$；这个条件意味着由 $\sigma$ 生成的规范对称保持~\eqref{eq:loopme} 的给定解不变。
此外，当 $t\not=\pm i$ 时，把方程 $\delta_T\eta=0$ 与 $\delta_T\tilde\eta=0$ 作一个线性组合便给出
% 式 (3.31)
\setcounter{equation}{30}
\begin{equation}\label{eq:woopem}
[\sigma,\bar\sigma]=0.
\end{equation}

只要我们只考虑不可约的超对称场——或者更一般地，考虑自同构群有限的超对称场——式~\eqref{eq:toopme} 就推出 $\sigma=0$。本文一般只考虑不可约的超对称场。关于可约情形的若干简单注记将出现在 \cite{ref26} 中。

在 $t=\pm i$ 处结论类似，只是细节不同。此时式~\eqref{eq:toopme} 说的是：$\sigma$ 生成平坦 $G_\Bbb{C}$ 值联络 $\CA$ 或 $\bar\CA$ 的一个对称。同样，只要只考虑不可约的超对称场，就可以假设 $\sigma=0$。（对 $t=\pm i$，若存在带 $\sigma\not=0$ 的可约超对称位形，我们得不到分离的条件 $D^*\phi=[\sigma,\bar\sigma]=0$，而只有线性组合 $D^*\phi\pm i[\sigma,\bar\sigma]=0$。）

对实 $t$，这里讨论的理论属于上同调场论 \cite{ref59} 的一般框架，就像四维的 Donaldson 理论或二维的 $A$-模型。场是具有 ${\EUK}=0$ 的玻色场 $A,\phi$；方程是~\eqref{eq:loopme}；对称性就是规范对称性。与任何这类理论一样，在没有算子插入时，路径积分可以通过对方程的解带符号计数来计算，只要这些解是孤立且不可约的。一个给定解的贡献符号，就是相应的单圈费米子行列式的符号。由空间 ${\eurm Y}_0$（它是 ${\eurm Y}$ 的一个分支）参数化的一族光滑紧的不可约解，其贡献是 ${\eurm Y}_0$ 的 Euler 示性数加上或减去一个符号，符号仍然来自费米子行列式。${\eurm Y}$ 中奇点的贡献以及可约解的贡献——例如 $A=\phi=0$ 的平凡解——更微妙，这里不予考虑。

由于带符号加权的解的数目在椭圆方程族的连续形变下不变，闭四维流形上无算子插入的配分函数必定与 $t$ 无关，至少对实 $t$ 如此。由全纯性，这对复 $t$ 也成立。闭四维流形上的配分函数与 $t$ 无关这一点，也可以更直接地从路径积分看出，后文会解释。从数学观点看，第~\ref{sec:vanishing}~节的消灭定理将给出关于 $t$-无关性的强陈述：式~\eqref{eq:loopme} 的解空间实际上与 $t$ 无关，$t=0,\infty$ 除外。

参数 $t$ 的重要性在于：可以引入一些算子，或者当 $M$ 有边界时引入一些边界条件，它们只在某个特定的 $t$ 值下保持拓扑对称性。与 Donaldson 理论一样，这里考虑的 TQFT 具有保持拓扑对称性的局部算子。这些算子在 \cite{ref53,ref54,ref55,ref56,ref57} 中讨论过，并将另文进一步分析 \cite{ref26}。然而，它们并不是几何 Langlands 纲领中最重要的算子。最重要的算子是第~\ref{sec:loopop}~节引入的 Wilson 与 't Hooft 线算子，它们在更熟悉的上同调场论中没有明显的近亲。

如果我们在实 $t$ 时的理论多少类比于二维 $A$-模型，那么复 $t$ 时它对应于什么？$t=\pm i$ 处的理论类比于二维 $B$-模型。在描述从黎曼面 $\Sigma$ 到复流形 $X$ 的映射 $\Phi:\Sigma\to X$ 的 $B$-模型中，超对称场是 $\Sigma$ 到 $X$ 的常映射。两种情形之间明显的相似之处是：常映射或平坦的 $G_{\Bbb{C}}$ 值联络都相对容易用超定方程（而非椭圆方程）描述——前者是 $d\Phi=0$，后者是 ${\cal F}=0$——而其中的复结构是显然的。当我们于第~\ref{sec:compacthitchin}~节考虑维数约化时，将会看到更深刻的类比。

\begin{yizhu}[物理图像]
按：把实 $t$ 与 $t=\pm i$ 两端对照起来看，本节给出了一个很有用的三分法：$t=\pm i$ 时方程是复的、超定的，解是平坦 $G_{\Bbb{C}}$ 值联络（类比二维 $B$-模型）；实 $t$ 时方程是实椭圆型的，路径积分按解的带符号计数给出（类比二维 $A$-模型，或四维 Donaldson 理论）；而 $t=0,\infty$ 是两个端点，方程退化，但通过乘以 $t$ 或 $t^{-1}$ 仍有意义（见式~\eqref{eq:loopme} 之后的讨论）。原文只谨慎地说这样得到的理论"松散地说是拓扑场论"，因为它可能依赖 $M$ 的光滑结构；但无论如何，$t$ 是这个理论族中唯一本质的参数（除 $t\sim-t$ 外），这也是为什么 $S$-对偶作用在 $t$ 上能产生非平凡的对偶性。
\end{yizhu}

\subsection{消灭定理}

\label{sec:vanishing}

令 $\EUV^+(t)= (F-\phi\wedge\phi+tD\phi)^+$、$\EUV^-(t)= (F-\phi\wedge\phi-t^{-1}D\phi)^-$、$\EUV_0= D^* \phi$，则超对称方程是 $\EUV^+(t)=\EUV^-(t)=\EUV_0=0$。本节我们集中讨论 $t$ 为实数的情形，此时超对称方程是椭圆型的。

与许多同超对称相联系的一阶方程一样——例如全纯曲线方程、Yang--Mills 瞬子方程、Seiberg--Witten 方程，或 \cite{ref45} 中研究的方程——这些方程受制于不寻常的消灭定理。在目前的情形，我们有：

\par\medskip
 {\it 消灭定理 1~}  设 $M$ 为紧的无边四维流形，$E$ 为 $M$ 上的 $G$-丛，其 Pontryagin 类非零，即 $\int_M\Tr\,F\wedge F\not=0$。那么对任何 $t\not=0,\infty$，超对称方程 $\EUV^+(t)=\EUV^-(t)=\EUV_0=0$ 无解。

\par\medskip  {\it 消灭定理 2~} 在这样的 $M$ 上，任何对某个不为 $0,\infty$ 的 $t$ 值满足这些方程的场，对一切 $t$ 都满足这些方程，因而由一个平坦的 $G_\Bbb{C}$ 值联络给出。

\begin{yizhu}[易错点]
按：消灭定理 1 是关于固定丛 $E$ 上的经典方程无解的陈述，而不是关于配分函数的陈述：它说的是对任意 $t\not=0,\infty$，方程 $\EUV^\pm(t)=\EUV_0=0$ 在 $E$ 上根本没有解，因此配分函数为零。这类"一阶方程 + 非零拓扑荷"的消灭现象在超对称方程中很常见（全纯曲线、Yang--Mills 瞬子、Seiberg--Witten 方程都有类似结论），其共同机制是把方程写成平方和等于拓扑荷的形式（见式~\eqref{eq:hombo} 与式~\eqref{eq:ormob}）：若平方和为零而拓扑荷非零，就不可能。定理 1 的证明只在 $t\not=0,\infty$ 有效，而 $t=0,\infty$ 处确实有解，这一点下文会说明。
\end{yizhu}

\par\medskip

第一个消灭定理的最短证明来自如下恒等式： % 式 (3.32)
\setcounter{equation}{31}
\begin{equation}\label{eq:hombo}
\int_M\Tr\,F\wedge
F=\int_M\Tr\left(\EUV^+(t)\wedge \EUV^+(-t^{-1})+\EUV^-(t)\wedge
\EUV^-(-t^{-1})\right).
\end{equation}
这个恒等式的证明是直接的，用到分部积分。显然，这个恒等式说明：若对某个 $t$ 有 $\EUV^\pm(t)=0$，则 $\int_M\Tr\, F\wedge F=0$。

当 $t=0,\infty$ 时，由于 $\EUV^\pm$ 的极点，这个证明失效。事实上，$t=0,\infty$ 时结论确实不成立：$t=0$ 或 $t=\infty$ 的方程至少有明显的解，即取 $\phi=0$ 的瞬子或反瞬子。

对第二消灭定理，我们需要一个类似但更复杂的恒等式： % 式 (3.33)
\setcounter{equation}{32}
\begin{equation}\label{eq:ormob}
\begin{aligned}&-\int_Md^4x\sqrt
g\,\Tr\,\left({t^{-1}\over
t+t^{-1}}\EUV^+(t)_{\mu\nu}\EUV^+(t)^{\mu\nu}+{t\over
t+t^{-1}}\EUV^-(t)_{\mu\nu}\EUV^-(t)^{\mu\nu}+\EUV_0^2\right)\\
&=-\int d^4x\sqrt g\,\Tr\,\left({1\over 2}F_{\mu\nu}
F^{\mu\nu}+D_\mu\phi_\nu D^\mu\phi^\nu
+R_{\mu\nu}\phi^\mu\phi^\nu+ {1\over
2}[\phi_\mu,\phi_\nu]^2\right)\\&~~~~~~~+{t-t^{-1}\over
(t+t^{-1})}\int_M\Tr\, F\wedge F. \\\end{aligned}
\end{equation}
这个恒等式令人惊讶之处在于：除了一项正比于 $\int\,\Tr\,F\wedge F$ 的拓扑不变量之外，右边与 $t$ 无关。其证明类似于式~\eqref{eq:hombo} 的证明，但更繁复。$R_{\mu\nu}$ 是 $M$ 的 Ricci 张量，它出现在把 $\int_Md^4x\sqrt g D_\mu\phi_\nu D^\nu\phi^\mu$ 与 $\int d^4x\sqrt g (D_\mu\phi^\mu)^2$ 联系起来的分部积分中。与式~\eqref{eq:hombo} 不同，式~\eqref{eq:ormob} 在 $t\to 0,\infty$ 时有自然的极限：只需把 $t (\EUV^-)^2/(t+t^{-1})$ 换成 $\lim_{t\to 0}(t\EUV^-)^2$，或把 $t^{-1}(\EUV^+)^2/(t+t^{-1})$ 换成 $\lim_{t\to \infty}(t^{-1}\EUV^+)^2$。因此基于式~\eqref{eq:ormob} 的论证在 $t=0,\infty$ 处也成立。

\begin{yizhu}[含糊点拨]
按：式~\eqref{eq:ormob} 是消灭定理的技术核心，值得看清它的结构：左边是 $-\int\sqrt g\,\Tr(\text{三组平方项})$，右边是 $-\int\sqrt g\,\Tr(\text{与 $t$ 无关的"能量"})$ 加上 $f(t)\int\Tr F\wedge F$。也就是说，两边之差给出了一个正比于拓扑荷的项，这是"对 $t$ 的依赖只能来自拓扑项"这一事实的定量形式；正因如此，$t\to 0,\infty$ 的极限存在，论证可以覆盖端点。它本质上是 Bogomolny 型平方和恒等式的高维版本，与四维瞬子方程、二维 $A$-模型的能量恒等式同源。第~\ref{sec:toplag}~节会通过超对称作用量的 $Q$-恰当分解给出它的物理解释。
\end{yizhu}

在第~\ref{sec:toplag}~节，我们将理解式~\eqref{eq:ormob} 的物理意义。现在，我们只用它来完成消灭定理的证明。

为证明第二消灭定理，我们必须证明：若丛 $E$ 上存在一对 $A,\phi$，使得对某个不为 $0,\infty$ 的 $t$ 值有 $\EUV^\pm(t)=\EUV_0=0$，则这些方程对一切 $t$ 都成立。我们已经知道，在这样的解中 $\int_M\Tr\,F\wedge F=0$，因此式~\eqref{eq:ormob} 的右边与 $t$ 无关。于是若左边对某个 $t$ 值取零，则它对一切 $t$ 都取零。此外，左边是平方和，因此对给定的 $t$，它取零当且仅当 $\EUV^\pm(t)=\EUV_0=0$。再者，如我们已经注意到的，$\EUV^\pm(t)$ 与 $\EUV_0$ 对一切 $t$ 取零，等价于 $G_{\Bbb{C}}$ 值联络 ${\cal A}=A+i\phi$ 的曲率 ${\cal F}=d{\cal A}+{\cal A}\wedge{\cal A}$ 取零。这就完成了论证。

最后，我们可以用式~\eqref{eq:ormob} 证明第一消灭定理的一个稍强版本：

\par\medskip{\it 消灭定理 1$'$~} 当 $\int_M\Tr\,F\wedge F>0$ 时，超对称方程除 $t=0$ 外无解；当 $\int_M\Tr\,F\wedge F<0$ 时，除 $t=\infty$ 外无解。\par\medskip

令 ${\eurm B}(t)$ 记式~\eqref{eq:ormob} 的左边，并令 $f(t)=(t-t^{-1})/(t+t^{-1})$。对任意 $t,u$，我们有
% 式 (3.34)
\setcounter{equation}{33}
\begin{equation}\label{eq:toto}
{\eurm B}(u)-{\eurm B}(t)=(f(u)-f(t))\int_M\Tr \,F\wedge
F.
\end{equation}
若 $\EUV^\pm(t)=\EUV_0=0$，则 ${\eurm B}(t)=0$；又因函数 ${\eurm B}$ 是正半定的，式~\eqref{eq:toto} 的左边是正半定的。为使右边正半定，若 $\int_M\Tr\,F\wedge F>0$，$t$ 必须是函数 $f$ 的最小值点；若 $\int_M\Tr\,F\wedge F<0$，$t$ 必须是最大值点。而唯一的最小值在 $t=0$，唯一的最大值在 $t=\infty$。

\begin{yizhu}[原文疑误]
按：这里有一个符号疑点，值得记录。式~\eqref{eq:ormob} 的左边显含一个负号，括号内的量对 $t\not=0$ 是正的平方和（系数的正负可直接验证），所以按字面计算 ${\eurm B}(t)\le 0$，即 ${\eurm B}$ 是负半定的；但本段接着说 ${\eurm B}$ 是正半定，并用"正半定"推理两次，最后得出 $t=0$ 为最小值点——定理 1$'$ 的陈述也与此一致。已核对的 arXiv 源文与已出版版本（PDF 第 31 页）此处均写作 positive semidefinite，故这是原文自身的符号不一致，而非抄录错误：要么式~\eqref{eq:ormob} 左边的负号应当去掉，要么"正半定/最小/最大"这串推理应整体取反。两种改法给出的 $\int\Tr F\wedge F$ 符号与 $t=0,\infty$ 的对应关系并不相同，读者若要引用该结论，宜先按自己的符号约定把式~\eqref{eq:ormob} 重新验算一遍。本书照原文逐字译出，未在正文改动。
\end{yizhu}

\par\medskip{\it 关于紧化的一个注记}

消灭定理的一个特例与紧化有关。我们取 $M=\Sigma\times C$，其中 $\Sigma$ 与 $C$ 是紧黎曼面。在我们的应用中，$C$ 是研究几何 Langlands 纲领所在的黎曼面，而我们把 $\Sigma$ 想成比 $C$ 大得多。不过在拓扑场论中，$\Sigma$ 与 $C$ 上的度规都无关紧要。

对 $t\not=0,\infty$，第二消灭定理说 $M$ 上超对称方程的任何解都由一个平坦的 $G_{\Bbb{C}}$ 值联络给出。$M=\Sigma\times C$ 上的平坦丛 $E$ 有分解\footnote{在这段讨论中，为简单起见我们使用向量丛的语言，就像 $G=SU(N)$ 一样。要把论证推广到任意 $G$，只需改用 $G$ 的子群而不是 $E$ 的子丛来表述即可。} $E=\oplus_{i=1}^n E'_i\otimes E''_i$，其中 $E'_i$ 是从 $\Sigma$ 拉回的平坦丛，$E''_i$ 是从 $C$ 拉回的平坦丛。

\begin{yizhu}[物理图像]
按：这个分解是本节与后文（第~\ref{sec:compacthitchin}~节）之间的桥梁，值得记住它的几何含义：$M=\Sigma\times C$ 上的平坦丛可以拆成"沿 $\Sigma$ 方向"与"沿 $C$ 方向"两部分的直和；限制到一点 $x\in\Sigma$ 之后，剩下的只是 $C$ 上的平坦丛 $E_x$，它正是 Hitchin 模空间 ${\MH}$ 中的一个点。于是 $\Sigma\times C$ 上的超对称场给出映射 $\Sigma\to{\MH}$；在 ${\EUN}=4$ 与 $S$-对偶的图景中，这正是四维理论紧化到二维后靶空间为 Hitchin 模空间的 sigma 模型的开端。Witten 2008 综述在讨论 Hitchin 纤维化与镜像对称时用的也是这套语言。
\end{yizhu}

若我们限制到一点 $x\in \Sigma$，则 $E'_i$ 变为平凡的。若 $d_i$ 是 $E_i'$ 的秩，则 $E$ 约化为 $E_x=\oplus_{i=1}^n d_iE''_i$，即从 $C$ 拉回的平坦丛的直和。$E_x$ 是 $C$ 上的平坦丛，因此定义了 Hitchin 模空间 ${\MH}$ 中的一个点；我们将在第~\ref{sec:compacthitchin}~节更充分地描述该模空间。眼下，我们可以就把 ${\MH}$ 想成 $C$ 上平坦 $G_\Bbb{C}$-丛的模空间。于是 $\Sigma\times C$ 上的任何超对称场都由一个从 $\Sigma$ 到 ${\MH}$ 的映射给出，即 $x\to E_x$。

模空间 ${\MH}$ 有对应于可约平坦丛的奇点。在当前的构造中，当 $n>1$ 时丛 $E_x$ 是可约的。对 $n=1$，有 $E=E'\otimes E''$，其中两个因子分别是从 $\Sigma$ 与 $C$ 拉回的。除非 $E'$ 的秩为 $1$（此时对 $G=SU(N)$ 它必为平凡的），$E_x$ 仍是可约的。因此，避开 ${\MH}$ 奇点的情形，恰好是 $E$ 为某个不可约平坦丛 $E''\to C$ 的拉回的情形。

超对称场同时也是使作用量取极小的位形；这一点在我们于第~\ref{sec:toplag}~节构造超对称拉格朗日量时就会清楚。在上述推理中，我们取 $\Sigma$ 与 $C$ 为紧的。在我们的应用里，我们也关心 $\Sigma$ 为完备但非紧的二维流形的情形，例如 $\Bbb{R}^2$、$\Bbb{R}^2_+$ 或 $\Bbb{R}\times I$（其中 $\Bbb{R}^2_+$ 是半空间，$I$ 是闭区间）。只要给定适当的边界条件与无穷远处的渐近行为，使得第二消灭定理的证明得以成立，这些情形下超对称场仍有分解 $E=\oplus_{i=1}^nE'_i\otimes E''_i$。（当 $\Sigma$ 平坦且单连通时，我们所说的 $E'_i$ 指带有平凡联络与协变常值 Higgs 场的丛。）由此得到与紧 $\Sigma$ 时相同的结论：避开 ${\MH}$ 奇点的超对称场是来自 $C$ 的拉回。
